\documentclass[11pt]{article}
\usepackage[a4paper,margin=1in]{geometry}
\usepackage{fontspec}
\usepackage[boldfont=false]{xeCJK}
\setCJKmainfont{FandolSong-Regular.otf}
\usepackage{amsmath}
\usepackage{amssymb}
\usepackage{array}
\usepackage{booktabs}
\usepackage{graphicx}
\usepackage{adjustbox}
\usepackage{enumitem}
\setlist[itemize]{topsep=2pt,partopsep=0pt,parsep=0pt,itemsep=2pt,leftmargin=1.4em}
\setlist[enumerate]{topsep=2pt,partopsep=0pt,parsep=0pt,itemsep=2pt,leftmargin=1.6em}
\usepackage{parskip}
\usepackage{fvextra}
\fvset{breaklines=true,breakanywhere=true,fontsize=\small}
\setlength{\parindent}{0pt}

\begin{document}

\begin{center}
  {\LARGE\bfseries Metals}\\[0.35em]
  {\large Igcse Chemistry}
\end{center}
\vspace{0.6em}

\section*{Properties of metals}

\subsection*{Physical properties}

\textbf{Metals} 金属 and \textbf{non-metals} 非金属 behave very differently:

\begin{center}
\adjustbox{max width=\linewidth}{%
\begin{tabular}{lll}
\toprule
\textbf{Property} & \textbf{Metals} & \textbf{Non-metals} \\
\midrule
\textbf{thermal conductivity} 导热性 (conducting heat) & good & poor \\
\textbf{electrical conductivity} 导电性 & good & poor (except graphite) \\
\textbf{malleability} 展性 and \textbf{ductility} 延性 & malleable and ductile & \textbf{brittle} 易碎 (they snap) \\
\textbf{melting point} 熔点 and \textbf{boiling point} 沸点 & usually high & usually low \\
\bottomrule
\end{tabular}}
\end{center}

\subsection*{Chemical properties}

Metals react in three main ways:

\begin{itemize}
\item with dilute \textbf{acids} 酸 → a \textbf{salt} 盐 + \textbf{hydrogen} 氢气

\item with cold water or \textbf{steam} 蒸汽 → a metal hydroxide (or oxide) + hydrogen

\item with \textbf{oxygen} 氧气 → a metal \textbf{oxide} 氧化物

\end{itemize}

\section*{Uses of metals}

A metal is chosen for a job because of its physical properties.

\begin{itemize}
\item \textbf{Aluminium} 铝 is used to make aircraft because of its low \textbf{density} 密度 (it is light).

\item Aluminium is used for overhead electrical cables because of its low density and good electrical conductivity.

\item Aluminium is used for food containers because it resists \textbf{corrosion} 腐蚀.

\item \textbf{Copper} 铜 is used for electrical wiring because of its good electrical conductivity and its ductility (it can be drawn into wires).

\end{itemize}

\section*{Alloys}

An \textbf{alloy} 合金 is a \textbf{mixture} 混合物 of a metal with one or more other elements.

\begin{itemize}
\item \textbf{Brass} 黄铜 is a mixture of copper and \textbf{zinc} 锌.

\item \textbf{Stainless steel} 不锈钢 is a mixture of \textbf{iron} 铁 with other elements such as \textbf{chromium} 铬, \textbf{nickel} 镍 and \textbf{carbon} 碳.

\end{itemize}

Alloys are usually harder and stronger than the pure metals, which makes them more useful. For example, stainless steel is used for cutlery because it is hard and does not rust.

\textbf{Why alloys are harder.} In a pure metal, the \textbf{atoms} 原子 are all the same size, so the \textbf{layers} 层 can slide over each other easily. In an alloy, the different-sized atoms stop the layers from sliding, so the alloy is harder and stronger.

\section*{The reactivity series}

The \textbf{reactivity series} 金属活动性顺序 lists metals in order of how reactive they are. Carbon and hydrogen are included for comparison:

\begin{quote}
\textbf{potassium} 钾, \textbf{sodium} 钠, \textbf{calcium} 钙, \textbf{magnesium} 镁, aluminium, carbon, zinc, iron, hydrogen, copper, silver, gold

\end{quote}

(most reactive at the top, least reactive at the bottom)

The higher a metal is, the more easily it forms positive \textbf{ions} 离子. This explains its reactions:

\begin{itemize}
\item potassium, sodium and calcium react with cold water.

\item magnesium reacts with steam (but only very slowly with cold water).

\item magnesium, zinc and iron react with dilute hydrochloric acid; copper, silver and gold do not.

\end{itemize}

\subsection*{Displacement reactions}

A more reactive metal will displace a less reactive metal from a solution of its ions (a \textbf{displacement reaction} 置换反应). For example, zinc displaces copper from copper(II) sulfate solution, because zinc is more reactive than copper.

\subsection*{The special case of aluminium}

Aluminium seems less reactive than its position suggests. This is because it is covered by a thin, strong \textbf{oxide layer} 氧化层 that stops other substances reaching the metal underneath.

\section*{Corrosion of metals}

\textbf{Rusting} 生锈 is the corrosion of iron and steel. Two things are needed for rusting: \textbf{oxygen} (from the air) and water. The \textbf{rust} 铁锈 formed is hydrated iron(III) oxide.

\subsection*{Stopping rust}

\textbf{Barrier methods} 隔离法 keep oxygen and water away from the iron:

\begin{itemize}
\item painting, \textbf{greasing} (covering with oil), and coating with plastic.

\end{itemize}

\textbf{Galvanising} 镀锌 means coating iron with a layer of zinc. This works in two ways:

\begin{itemize}
\item it is a barrier (the zinc keeps out air and water);

\item it gives \textbf{sacrificial protection} 牺牲保护. Because zinc is more reactive than iron, the zinc loses \textbf{electrons} 电子 and corrodes instead of the iron — even if the surface is scratched.

\end{itemize}

\section*{Extraction of metals}

How a metal is taken from its \textbf{ore} 矿石 depends on its place in the reactivity series. Metals below carbon can be extracted by heating with carbon (which removes the oxygen). Metals above carbon are too reactive for this and must be extracted by \textbf{electrolysis} 电解.

\subsection*{Iron from the blast furnace}

Iron is extracted from its ore \textbf{hematite} 赤铁矿 (iron(III) oxide) in a \textbf{blast furnace} 高炉:

\begin{itemize}
\item Carbon (coke) burns in hot air to give carbon dioxide and lots of heat.

\item This \textbf{carbon dioxide} 二氧化碳 reacts with more carbon to form \textbf{carbon monoxide} 一氧化碳.

\item The carbon monoxide \textbf{reduces} the iron(III) oxide to iron (this is \textbf{reduction} 还原).

\item \textbf{Limestone} 石灰石 (calcium carbonate) breaks down in the heat to form \textbf{calcium oxide} 氧化钙.

\item The calcium oxide reacts with sandy impurities to form \textbf{slag} 炉渣, which is removed.

\end{itemize}

\subsection*{Aluminium by electrolysis}

Aluminium is extracted from its ore \textbf{bauxite} 铝土矿 (purified to aluminium oxide) by electrolysis:

\begin{itemize}
\item The aluminium oxide is dissolved in molten \textbf{cryolite} 冰晶石 to lower its melting point and save energy.

\item At the \textbf{cathode} 阴极, aluminium ions gain electrons to form aluminium metal.

\item At the \textbf{anode} 阳极, oxygen is formed. This oxygen reacts with the hot carbon anodes and burns them away, so they must be replaced regularly.

\end{itemize}


\end{document}
